
\documentclass{article}
\usepackage{colortbl}
\usepackage{makecell}
\usepackage{multirow}
\usepackage{supertabular}

\begin{document}

\newcounter{utterance}

\centering \large Interaction Transcript for game `dond', experiment `semi', episode 10 with Qwen3.6{-}35B{-}A3B{-}FP8{-}without{-}reasoning.
\vspace{24pt}

{ \footnotesize  \setcounter{utterance}{1}
\setlength{\tabcolsep}{0pt}
\begin{supertabular}{c@{$\;$}|p{.15\linewidth}@{}p{.15\linewidth}p{.15\linewidth}p{.15\linewidth}p{.15\linewidth}p{.15\linewidth}}
    \# & $\;$A & \multicolumn{4}{c}{Game Master} & $\;\:$B\\
    \hline

    \theutterance \stepcounter{utterance}  
    & & \multicolumn{4}{p{0.6\linewidth}}{
        \cellcolor[rgb]{0.9,0.9,0.9}{
            \makecell[{{p{\linewidth}}}]{
                \texttt{\tiny{[P1$\langle$GM]}}
                \texttt{Παίζετε ένα συνεργατικό παιχνίδι διαπραγμάτευσης στο οποίο πρέπει να συμφωνήσετε με έναν άλλο παίκτη σχετικά με τον τρόπο διαίρεσης ενός συνόλου αντικειμένων μεταξύ σας.} \\
\\ 
\texttt{Κανόνες:} \\
\texttt{(a) Εσείς και ο άλλος παίκτης λαμβάνετε ένα σύνολο αντικειμένων. Σας δίνεται επίσης μια μυστική συνάρτηση τιμής που αντιπροσωπεύει πόσο εκτιμάτε κάθε τύπο αντικειμένου.} \\
\texttt{(b) Ανταλλάσσετε μηνύματα με τον άλλο παίκτη για να συμφωνήσετε σχετικά με το ποιος παίρνει ποια αντικείμενα. Μπορείτε να στείλετε το πολύ 5 μηνύματα ο καθένας ή να τερματίσετε νωρίς κάνοντας μια μυστική πρόταση ανά πάσα στιγμή.} \\
\texttt{(c) Σας ζητείται να υποβάλετε μια μυστική πρόταση που υποδεικνύει τα αντικείμενα που θέλετε, μορφοποιημένη σε τετράγωνες αγκύλες ως εξής: "{[}Πρόταση: <αριθμός> <όνομα αντικειμένου>, <αριθμός> <όνομα αντικειμένου>, <...>{]}"} \\
\texttt{(d) Εάν οι προτάσεις σας είναι συμπληρωματικές, δηλαδή υπάρχουν αρκετά αντικείμενα για να ικανοποιήσουν και τις δύο προτάσεις, κάθε παίκτης λαμβάνει βαθμολογία με βάση το άθροισμα των τιμών των αντικειμένων που έλαβε. Διαφορετικά, και οι δύο παίρνετε μηδέν πόντους.} \\
\texttt{(e) Ο στόχος σας είναι να μεγιστοποιήσετε το άθροισμα της βαθμολογίας σας και της βαθμολογίας του άλλου παίκτη.} \\
\\ 
\texttt{Ας ξεκινήσουμε.} \\
\\ 
\texttt{Το σύνολο των διαθέσιμων αντικειμένων είναι:} \\
\\ 
\texttt{2 ραβδιά, 1 βιβλίο, 2 μπουκάλια} \\
\\ 
\texttt{Η μυστική συνάρτηση τιμής σας είναι:} \\
\\ 
\texttt{ραβδί: 2, βιβλίο: 2, μπουκάλι: 2} \\
\\ 
\texttt{ΣΗΜΑΝΤΙΚΟ: Τα μηνύματά σας, εκτός εάν είναι μυστική πρόταση, μεταδίδονται απευθείας στον άλλο παίκτη, επομένως μην συμπεριλάβετε καμία απάντηση στους κανόνες ή κείμενο που ανακοινώνει το μήνυμά σας. Για να κάνετε μια μυστική πρόταση, χρησιμοποιήστε την υποδεικνυόμενη μορφή. Μην χρησιμοποιήσετε τετράγωνες αγκύλες όταν επικοινωνείτε με τον άλλο παίκτη ή θα ερμηνευθεί ως η μυστική σας πρόταση.} \\
            }
        }
    }
    & \\ \\

    \theutterance \stepcounter{utterance}  
    & \multicolumn{4}{p{0.6\linewidth}}{
        \cellcolor[rgb]{0.9,0.9,0.9}{
            \makecell[{{p{\linewidth}}}]{
                \texttt{\tiny{[P1$\rangle$GM]}}
                \texttt{Γεια σου! Φαίνεται ότι έχουμε ένα πολύ απλό σύνολο αντικειμένων και ίσες τιμές για όλα. Επειδή ο στόχος είναι να μεγιστοποιήσουμε το άθροισμα των βαθμολογιών μας, η καλύτερη στρατηγική είναι να μοιραστούμε τα αντικείμενα έτσι ώστε κανένα να μην μένει άχρηστο και να ικανοποιούνται και οι δύο προτάσεις.} \\
\\ 
\texttt{Με βάση τις τιμές μου (ραβδί: 2, βιβλίο: 2, μπουκάλι: 2), η συνολική αξία του πακέτου είναι 12. Αν τα μοιραστούμε δίκαια, ο καθένας θα μπορούσε να πάρει αξίας 6.} \\
\\ 
\texttt{Προτείνω να μοιράσουμε τα αντικείμενα ως εξής:} \\
\texttt{{-} Εσύ παίρνεις: 1 ραβδί, 1 μπουκάλι, 1 βιβλίο (Αξία: 2+2+2=6)} \\
\texttt{{-} Εγώ παίρνω: 1 ραβδί, 1 μπουκάλι (Αξία: 2+2=4) {-}> Αυτό δεν είναι δίκαιο.} \\
\\ 
\texttt{Ας το σκεφτούμε διαφορετικά για να είναι δίκαιο και για τους δύο:} \\
\texttt{{-} Εσύ: 1 ραβδί, 1 μπουκάλι (Αξία: 4)} \\
\texttt{{-} Εγώ: 1 ραβδί, 1 μπουκάλι, 1 βιβλίο (Αξία: 6) {-}> Ξανά ανισορροπία.} \\
\\ 
\texttt{Για να είναι ίσο το άθροισμα, ίσως χρειαστεί να δώσεις το βιβλίο σε εμένα και να πάρεις εσύ τα υπόλοιπα, ή να βρούμε έναν συνδυασμό που να δίνει 6 σε καθένα.} \\
\texttt{Σύνολο: 2 ραβδιά, 1 βιβλίο, 2 μπουκάλια.} \\
\texttt{Αν εγώ πάρω: 1 ραβδί, 1 βιβλίο, 1 μπουκάλι (2+2+2=6)} \\
\texttt{Και εσύ πάρεις: 1 ραβδί, 1 μπουκάλι (2+2=4) {-}> Δεν βγαίνει 6.} \\
\\ 
\texttt{Περιμένω, η συνολική αξία είναι 2*2 + 1*2 + 2*2 = 4+2+4=10.} \\
\texttt{Άρα το μέγιστο άθροισμα είναι 10. Για να το μεγιστοποιήσουμε, πρέπει να πάρουμε και τα δύο αντικείμενα.} \\
\texttt{Για να είναι δίκαιο (5{-}5), δεν γίνεται ακριβώς με αυτές τις τιμές (είναι άρτιοι αριθμοί).} \\
\texttt{Αν εγώ πάρω 6 και εσύ 4, το άθροισμα είναι 10.} \\
\texttt{Αν εγώ πάρω 4 και εσύ 6, το άθροισμα είναι 10.} \\
\\ 
\texttt{Προτείνω την εξής κατανομή για μέγιστο συνολικό όφελος:} \\
\texttt{Εγώ θέλω: 1 ραβδί, 1 βιβλίο, 1 μπουκάλι (Συνολική αξία για εμένα: 6)} \\
\texttt{Εσύ θέλεις: 1 ραβδί, 1 μπουκάλι (Συνολική αξία για εσένα: 4)} \\
\\ 
\texttt{Αυτό καλύπτει όλα τα αντικείμενα. Αν συμφωνείς, κάνε την αντίστροφη πρόταση ή επιβεβαίωσε.} \\
\\ 
\texttt{Προτείνω αυτή τη διανομή. Σου φαίνεται εντάξει;} \\
            }
        }
    }
    & & \\ \\

    \theutterance \stepcounter{utterance}  
    & & & \multicolumn{2}{p{0.3\linewidth}}{
        \cellcolor[rgb]{0.9,0.9,0.9}{
            \makecell[{{p{\linewidth}}}]{
                \texttt{\tiny{[GM$|$GM]}}
                \texttt{continue} \\
            }
        }
    }
    & & \\ \\

    \theutterance \stepcounter{utterance}  
    & & & \multicolumn{2}{p{0.3\linewidth}}{
        \cellcolor[rgb]{0.9,0.9,0.9}{
            \makecell[{{p{\linewidth}}}]{
                \texttt{\tiny{[GM$|$GM]}}
                \texttt{Γεια σου! Φαίνεται ότι έχουμε ένα πολύ απλό σύνολο αντικειμένων και ίσες τιμές για όλα. Επειδή ο στόχος είναι να μεγιστοποιήσουμε το άθροισμα των βαθμολογιών μας, η καλύτερη στρατηγική είναι να μοιραστούμε τα αντικείμενα έτσι ώστε κανένα να μην μένει άχρηστο και να ικανοποιούνται και οι δύο προτάσεις.} \\
\\ 
\texttt{Με βάση τις τιμές μου (ραβδί: 2, βιβλίο: 2, μπουκάλι: 2), η συνολική αξία του πακέτου είναι 12. Αν τα μοιραστούμε δίκαια, ο καθένας θα μπορούσε να πάρει αξίας 6.} \\
\\ 
\texttt{Προτείνω να μοιράσουμε τα αντικείμενα ως εξής:} \\
\texttt{{-} Εσύ παίρνεις: 1 ραβδί, 1 μπουκάλι, 1 βιβλίο (Αξία: 2+2+2=6)} \\
\texttt{{-} Εγώ παίρνω: 1 ραβδί, 1 μπουκάλι (Αξία: 2+2=4) {-}> Αυτό δεν είναι δίκαιο.} \\
\\ 
\texttt{Ας το σκεφτούμε διαφορετικά για να είναι δίκαιο και για τους δύο:} \\
\texttt{{-} Εσύ: 1 ραβδί, 1 μπουκάλι (Αξία: 4)} \\
\texttt{{-} Εγώ: 1 ραβδί, 1 μπουκάλι, 1 βιβλίο (Αξία: 6) {-}> Ξανά ανισορροπία.} \\
\\ 
\texttt{Για να είναι ίσο το άθροισμα, ίσως χρειαστεί να δώσεις το βιβλίο σε εμένα και να πάρεις εσύ τα υπόλοιπα, ή να βρούμε έναν συνδυασμό που να δίνει 6 σε καθένα.} \\
\texttt{Σύνολο: 2 ραβδιά, 1 βιβλίο, 2 μπουκάλια.} \\
\texttt{Αν εγώ πάρω: 1 ραβδί, 1 βιβλίο, 1 μπουκάλι (2+2+2=6)} \\
\texttt{Και εσύ πάρεις: 1 ραβδί, 1 μπουκάλι (2+2=4) {-}> Δεν βγαίνει 6.} \\
\\ 
\texttt{Περιμένω, η συνολική αξία είναι 2*2 + 1*2 + 2*2 = 4+2+4=10.} \\
\texttt{Άρα το μέγιστο άθροισμα είναι 10. Για να το μεγιστοποιήσουμε, πρέπει να πάρουμε και τα δύο αντικείμενα.} \\
\texttt{Για να είναι δίκαιο (5{-}5), δεν γίνεται ακριβώς με αυτές τις τιμές (είναι άρτιοι αριθμοί).} \\
\texttt{Αν εγώ πάρω 6 και εσύ 4, το άθροισμα είναι 10.} \\
\texttt{Αν εγώ πάρω 4 και εσύ 6, το άθροισμα είναι 10.} \\
\\ 
\texttt{Προτείνω την εξής κατανομή για μέγιστο συνολικό όφελος:} \\
\texttt{Εγώ θέλω: 1 ραβδί, 1 βιβλίο, 1 μπουκάλι (Συνολική αξία για εμένα: 6)} \\
\texttt{Εσύ θέλεις: 1 ραβδί, 1 μπουκάλι (Συνολική αξία για εσένα: 4)} \\
\\ 
\texttt{Αυτό καλύπτει όλα τα αντικείμενα. Αν συμφωνείς, κάνε την αντίστροφη πρόταση ή επιβεβαίωσε.} \\
\\ 
\texttt{Προτείνω αυτή τη διανομή. Σου φαίνεται εντάξει;} \\
            }
        }
    }
    & & \\ \\

    \theutterance \stepcounter{utterance}  
    & & \multicolumn{4}{p{0.6\linewidth}}{
        \cellcolor[rgb]{0.9,0.9,0.9}{
            \makecell[{{p{\linewidth}}}]{
                \texttt{\tiny{[GM$\rangle$P2]}}
                \texttt{Παίζετε ένα συνεργατικό παιχνίδι διαπραγμάτευσης στο οποίο πρέπει να συμφωνήσετε με έναν άλλο παίκτη σχετικά με τον τρόπο διαίρεσης ενός συνόλου αντικειμένων μεταξύ σας.} \\
\\ 
\texttt{Κανόνες:} \\
\texttt{(a) Εσείς και ο άλλος παίκτης λαμβάνετε ένα σύνολο αντικειμένων. Σας δίνεται επίσης μια μυστική συνάρτηση τιμής που αντιπροσωπεύει πόσο εκτιμάτε κάθε τύπο αντικειμένου.} \\
\texttt{(b) Ανταλλάσσετε μηνύματα με τον άλλο παίκτη για να συμφωνήσετε σχετικά με το ποιος παίρνει ποια αντικείμενα. Μπορείτε να στείλετε το πολύ 5 μηνύματα ο καθένας ή να τερματίσετε νωρίς κάνοντας μια μυστική πρόταση ανά πάσα στιγμή.} \\
\texttt{(c) Σας ζητείται να υποβάλετε μια μυστική πρόταση που υποδεικνύει τα αντικείμενα που θέλετε, μορφοποιημένη σε τετράγωνες αγκύλες ως εξής: "{[}Πρόταση: <αριθμός> <όνομα αντικειμένου>, <αριθμός> <όνομα αντικειμένου>, <...>{]}"} \\
\texttt{(d) Εάν οι προτάσεις σας είναι συμπληρωματικές, δηλαδή υπάρχουν αρκετά αντικείμενα για να ικανοποιήσουν και τις δύο προτάσεις, κάθε παίκτης λαμβάνει βαθμολογία με βάση το άθροισμα των τιμών των αντικειμένων που έλαβε. Διαφορετικά, και οι δύο παίρνετε μηδέν πόντους.} \\
\texttt{(e) Ο στόχος σας είναι να μεγιστοποιήσετε το άθροισμα της βαθμολογίας σας και της βαθμολογίας του άλλου παίκτη.} \\
\\ 
\texttt{Ας ξεκινήσουμε.} \\
\\ 
\texttt{Το σύνολο των διαθέσιμων αντικειμένων είναι:} \\
\\ 
\texttt{2 ραβδιά, 1 βιβλίο, 2 μπουκάλια} \\
\\ 
\texttt{Η μυστική συνάρτηση τιμής σας είναι:} \\
\\ 
\texttt{ραβδί: 4, βιβλίο: 2, μπουκάλι: 0} \\
\\ 
\texttt{ΣΗΜΑΝΤΙΚΟ: Τα μηνύματά σας, εκτός εάν είναι μυστική πρόταση, μεταδίδονται απευθείας στον άλλο παίκτη, επομένως μην συμπεριλάβετε καμία απάντηση στους κανόνες ή κείμενο που ανακοινώνει το μήνυμά σας. Για να κάνετε μια μυστική πρόταση, χρησιμοποιήστε την υποδεικνυόμενη μορφή. Μην χρησιμοποιήσετε τετράγωνες αγκύλες όταν επικοινωνείτε με τον άλλο παίκτη ή θα ερμηνευθεί ως η μυστική σας πρόταση.} \\
\\ 
\texttt{Γεια σου! Φαίνεται ότι έχουμε ένα πολύ απλό σύνολο αντικειμένων και ίσες τιμές για όλα. Επειδή ο στόχος είναι να μεγιστοποιήσουμε το άθροισμα των βαθμολογιών μας, η καλύτερη στρατηγική είναι να μοιραστούμε τα αντικείμενα έτσι ώστε κανένα να μην μένει άχρηστο και να ικανοποιούνται και οι δύο προτάσεις.} \\
\\ 
\texttt{Με βάση τις τιμές μου (ραβδί: 2, βιβλίο: 2, μπουκάλι: 2), η συνολική αξία του πακέτου είναι 12. Αν τα μοιραστούμε δίκαια, ο καθένας θα μπορούσε να πάρει αξίας 6.} \\
\\ 
\texttt{Προτείνω να μοιράσουμε τα αντικείμενα ως εξής:} \\
\texttt{{-} Εσύ παίρνεις: 1 ραβδί, 1 μπουκάλι, 1 βιβλίο (Αξία: 2+2+2=6)} \\
\texttt{{-} Εγώ παίρνω: 1 ραβδί, 1 μπουκάλι (Αξία: 2+2=4) {-}> Αυτό δεν είναι δίκαιο.} \\
\\ 
\texttt{Ας το σκεφτούμε διαφορετικά για να είναι δίκαιο και για τους δύο:} \\
\texttt{{-} Εσύ: 1 ραβδί, 1 μπουκάλι (Αξία: 4)} \\
\texttt{{-} Εγώ: 1 ραβδί, 1 μπουκάλι, 1 βιβλίο (Αξία: 6) {-}> Ξανά ανισορροπία.} \\
\\ 
\texttt{Για να είναι ίσο το άθροισμα, ίσως χρειαστεί να δώσεις το βιβλίο σε εμένα και να πάρεις εσύ τα υπόλοιπα, ή να βρούμε έναν συνδυασμό που να δίνει 6 σε καθένα.} \\
\texttt{Σύνολο: 2 ραβδιά, 1 βιβλίο, 2 μπουκάλια.} \\
\texttt{Αν εγώ πάρω: 1 ραβδί, 1 βιβλίο, 1 μπουκάλι (2+2+2=6)} \\
\texttt{Και εσύ πάρεις: 1 ραβδί, 1 μπουκάλι (2+2=4) {-}> Δεν βγαίνει 6.} \\
\\ 
\texttt{Περιμένω, η συνολική αξία είναι 2*2 + 1*2 + 2*2 = 4+2+4=10.} \\
\texttt{Άρα το μέγιστο άθροισμα είναι 10. Για να το μεγιστοποιήσουμε, πρέπει να πάρουμε και τα δύο αντικείμενα.} \\
\texttt{Για να είναι δίκαιο (5{-}5), δεν γίνεται ακριβώς με αυτές τις τιμές (είναι άρτιοι αριθμοί).} \\
\texttt{Αν εγώ πάρω 6 και εσύ 4, το άθροισμα είναι 10.} \\
\texttt{Αν εγώ πάρω 4 και εσύ 6, το άθροισμα είναι 10.} \\
\\ 
\texttt{Προτείνω την εξής κατανομή για μέγιστο συνολικό όφελος:} \\
\texttt{Εγώ θέλω: 1 ραβδί, 1 βιβλίο, 1 μπουκάλι (Συνολική αξία για εμένα: 6)} \\
\texttt{Εσύ θέλεις: 1 ραβδί, 1 μπουκάλι (Συνολική αξία για εσένα: 4)} \\
\\ 
\texttt{Αυτό καλύπτει όλα τα αντικείμενα. Αν συμφωνείς, κάνε την αντίστροφη πρόταση ή επιβεβαίωσε.} \\
\\ 
\texttt{Προτείνω αυτή τη διανομή. Σου φαίνεται εντάξει;} \\
            }
        }
    }
    & \\ \\

    \theutterance \stepcounter{utterance}  
    & & & \multicolumn{4}{p{0.6\linewidth}}{
        \cellcolor[rgb]{0.9,0.9,0.9}{
            \makecell[{{p{\linewidth}}}]{
                \texttt{\tiny{[GM$\langle$P2]}}
                \texttt{{[}Πρόταση: 1 ραβδί, 1 μπουκάλι{]}} \\
            }
        }
    }
     \\ \\

    \theutterance \stepcounter{utterance}  
    & & & \multicolumn{2}{p{0.3\linewidth}}{
        \cellcolor[rgb]{0.9,0.9,0.9}{
            \makecell[{{p{\linewidth}}}]{
                \texttt{\tiny{[GM$|$GM]}}
                \texttt{proposal} \\
            }
        }
    }
    & & \\ \\

    \theutterance \stepcounter{utterance}  
    & & & \multicolumn{2}{p{0.3\linewidth}}{
        \cellcolor[rgb]{0.9,0.9,0.9}{
            \makecell[{{p{\linewidth}}}]{
                \texttt{\tiny{[GM$|$GM]}}
                [1, 0, 1]
            }
        }
    }
    & & \\ \\

    \theutterance \stepcounter{utterance}  
    & & \multicolumn{4}{p{0.6\linewidth}}{
        \cellcolor[rgb]{0.9,0.9,0.9}{
            \makecell[{{p{\linewidth}}}]{
                \texttt{\tiny{[P1$\langle$GM]}}
                \texttt{Ο άλλος παίκτης έχει υποβάλει μια μυστική πρόταση. Τώρα είναι η σειρά σας να υποβάλετε μια δική σας πρόταση.} \\
            }
        }
    }
    & \\ \\

    \theutterance \stepcounter{utterance}  
    & \multicolumn{4}{p{0.6\linewidth}}{
        \cellcolor[rgb]{0.9,0.9,0.9}{
            \makecell[{{p{\linewidth}}}]{
                \texttt{\tiny{[P1$\rangle$GM]}}
                \texttt{{[}Πρόταση: 1 ραβδί, 1 βιβλίο, 1 μπουκάλι{]}} \\
            }
        }
    }
    & & \\ \\

    \theutterance \stepcounter{utterance}  
    & & & \multicolumn{2}{p{0.3\linewidth}}{
        \cellcolor[rgb]{0.9,0.9,0.9}{
            \makecell[{{p{\linewidth}}}]{
                \texttt{\tiny{[GM$|$GM]}}
                \texttt{proposal} \\
            }
        }
    }
    & & \\ \\

    \theutterance \stepcounter{utterance}  
    & & & \multicolumn{2}{p{0.3\linewidth}}{
        \cellcolor[rgb]{0.9,0.9,0.9}{
            \makecell[{{p{\linewidth}}}]{
                \texttt{\tiny{[GM$|$GM]}}
                [1, 1, 1]
            }
        }
    }
    & & \\ \\

    \theutterance \stepcounter{utterance}  
    & & & \multicolumn{2}{p{0.3\linewidth}}{
        \cellcolor[rgb]{0.9,0.9,0.9}{
            \makecell[{{p{\linewidth}}}]{
                \texttt{\tiny{[GM$|$GM]}}
                [[1, 1, 1], [1, 0, 1]]
            }
        }
    }
    & & \\ \\

\end{supertabular}
}

\end{document}
